\section{Síntesis y ampliación}

\subsection{Repaso de los puntos clave de esta clase}

La primera clase de MIT 6.S191 presentó de manera sistemática, de lo general a lo particular, el marco fundamental del aprendizaje profundo:

\begin{enumerate}
\item \textbf{Ubicación del aprendizaje profundo}: IA $\supset$ ML $\supset$ DL, cuyo núcleo consiste en aprender de manera automática representaciones de características a partir de datos crudos
\item \textbf{Perceptrón}: la unidad básica de las redes neuronales, que consta de tres pasos: suma ponderada, sesgo y activación no lineal
\item \textbf{Funciones de activación}: funciones no lineales como Sigmoid, Tanh y ReLU que permiten a la red aproximar funciones arbitrariamente complejas
\item \textbf{Redes profundas}: se construyen apilando múltiples capas Dense, y cada capa aprende características cada vez más abstractas
\item \textbf{Funciones de pérdida}: la entropía cruzada (clasificación) y el MSE (regresión) miden respectivamente la distancia entre la predicción y el valor real
\item \textbf{Proceso de entrenamiento}: descenso de gradiente + retropropagación + una tasa de aprendizaje adecuada
\item \textbf{Regularización}: Dropout y Early Stopping previenen el sobreajuste
\end{enumerate}

\begin{table}[H]
\centering
\begin{tabular}{lp{5cm}p{5cm}}
\toprule
\textbf{Concepto} & \textbf{Expresión matemática} & \textbf{Explicación intuitiva} \\
\midrule
Perceptrón & $\hat{y} = g(w_0 + \mathbf{X}^T\mathbf{W})$ & Voto ponderado + decisión no lineal \\
Capa Dense & $z_i = w_{0,i} + \sum_j x_j w_{j,i}$ & Varios perceptrones en paralelo \\
Función de pérdida & $J(\mathbf{W}) = \frac{1}{n}\sum_i \mathcal{L}$ & Medida cuantitativa de la calidad de la predicción \\
Descenso de gradiente & $\mathbf{W} \leftarrow \mathbf{W} - \eta \nabla J$ & Bajar la montaña en la dirección más pronunciada \\
Retropropagación & Aplicación recursiva de la regla de la cadena & Propaga el error capa por capa, de la salida hacia la entrada \\
Dropout & Probabilidad $p$ de anular aleatoriamente & Fuerza la redundancia y evita la dependencia de una sola característica \\
\bottomrule
\end{tabular}
\caption{Tabla de referencia rápida de los conceptos clave de esta clase}
\end{table}

\subsection{De lo fundamental a lo más avanzado}

Amini enfatizó varias veces durante la clase que el contenido de esta sesión (el perceptrón, el descenso de gradiente, la retropropagación) son invenciones de hace décadas, pero constituyen la base necesaria para entender el aprendizaje profundo moderno. Las siguientes clases del curso cubrirán:

\begin{itemize}
\item \textbf{Modelado de secuencias}: arquitecturas como RNN y LSTM para procesar datos temporales
\item \textbf{Visión por computadora}: CNN y arquitecturas de visión modernas
\item \textbf{Modelos generativos}: VAE, GAN, Diffusion Models
\item \textbf{Aprendizaje por refuerzo}: hacer que el modelo aprenda una política mediante la interacción con el entorno
\item \textbf{Modelos de lenguaje de gran escala}: la arquitectura Transformer y los LLM modernos
\end{itemize}

\subsection{Lecturas adicionales}

\begin{itemize}
\item Sitio oficial del curso: \url{https://introtodeeplearning.com}
\item Repositorio de código de los laboratorios: \url{https://github.com/MITDeepLearning/introtodeeplearning}
\item Ian Goodfellow, Yoshua Bengio, Aaron Courville, \textit{Deep Learning}, MIT Press, 2016
\item Michael Nielsen, \textit{Neural Networks and Deep Learning}, libro de texto gratuito en línea
\item 3Blue1Brown, serie de videos \textit{Neural Networks} (visualización intuitiva)
\end{itemize}

%% --- Fin del contenido --- %%

\end{document}
